% Options for packages loaded elsewhere
\PassOptionsToPackage{unicode}{hyperref}
\PassOptionsToPackage{hyphens}{url}
%
\documentclass[
]{article}
\usepackage{amsmath,amssymb}
\usepackage{lmodern}
\usepackage{iftex}
\ifPDFTeX
  \usepackage[T1]{fontenc}
  \usepackage[utf8]{inputenc}
  \usepackage{textcomp} % provide euro and other symbols
\else % if luatex or xetex
  \usepackage{unicode-math}
  \defaultfontfeatures{Scale=MatchLowercase}
  \defaultfontfeatures[\rmfamily]{Ligatures=TeX,Scale=1}
\fi
% Use upquote if available, for straight quotes in verbatim environments
\IfFileExists{upquote.sty}{\usepackage{upquote}}{}
\IfFileExists{microtype.sty}{% use microtype if available
  \usepackage[]{microtype}
  \UseMicrotypeSet[protrusion]{basicmath} % disable protrusion for tt fonts
}{}
\makeatletter
\@ifundefined{KOMAClassName}{% if non-KOMA class
  \IfFileExists{parskip.sty}{%
    \usepackage{parskip}
  }{% else
    \setlength{\parindent}{0pt}
    \setlength{\parskip}{6pt plus 2pt minus 1pt}}
}{% if KOMA class
  \KOMAoptions{parskip=half}}
\makeatother
\usepackage{xcolor}
\setlength{\emergencystretch}{3em} % prevent overfull lines
\providecommand{\tightlist}{%
  \setlength{\itemsep}{0pt}\setlength{\parskip}{0pt}}
\setcounter{secnumdepth}{-\maxdimen} % remove section numbering
\ifLuaTeX
  \usepackage{selnolig}  % disable illegal ligatures
\fi
\IfFileExists{bookmark.sty}{\usepackage{bookmark}}{\usepackage{hyperref}}
\IfFileExists{xurl.sty}{\usepackage{xurl}}{} % add URL line breaks if available
\urlstyle{same} % disable monospaced font for URLs
\hypersetup{
  hidelinks,
  pdfcreator={LaTeX via pandoc}}

\author{Dietmar Gerald Schrausser}
\date{09.05.2026}


\begin{document}

\hypertarget{foliumblat-iiiir}{%
\section{Foliu{[}m{]}/Blat IIIIr}\label{foliumblat-iiiir}}

Eine zeitgenössische deutsche Übersetzung wird präsentiert, die auf
einem Vergleich der bestehenden lateinischen, frühneuhochdeutschen
(ENHG) und englischen (Übersetzungs-) Texte basiert, um ein besseres
Verständnis zu ermöglichen. Insbesondere da der ENHG-Text im Vergleich
zum modernen Hochdeutschen schwer verständlich ist, da dieser fundierte
Deutschkenntnisse (als Muttersprache) sowie Kenntnisse verschiedener
Dialekte voraussetzt.

\hypertarget{vom-werk-des-vierten-tages}{%
\subsection{Vom Werk des vierten
Tages}\label{vom-werk-des-vierten-tages}}

\begin{quote}
\textbf{Q}Uarto die dixit deus. Fia{[}n{]}t luminaria in
firmame{[}n{]}to celi: {[}et{]} diuidant diem et nocte{[}m{]}.\\
Et sint in signa {[}et{]} t{[}em{]}p{[}er{]}a: {[}et{]} dies: {[}et{]}
a{[}n{]}nos: vt lucea{[}n{]}t in firmame{[}n{]}to celi: {[}et{]}
illumine{[}n{]}t terra{[}m{]}.\\
Et factu{[}m{]} e{[}st{]} ita.\\
Fecitq{[}ue{]} deus duos luminaria magna: luminare maius: vt
p{[}re{]}esset diei: {[}et{]} luminare min{[}us{]}: vt p{[}er{]}esset
nocti {[}et{]} stellas: vt duiderent luce{[}m{]} {[}et{]} tenebras.
\end{quote}

\textbf{Moses} erwähnt zuallererst die Himmelskörper\footnote{`celestium'
  und `himlischen ding'.}, die Gott in das Firmament gesetzt hat, damit
sie am Himmel\footnote{`celo' und `himel'.} strahlen und um die
\emph{Erde} zu erleuchten, also Sonne, Mond und Sterne\footnote{`stellas'
  und `stern'.} mit denen der obere\footnote{`superior' und `oberteil'.}
Teil der \emph{Welt} selbst geschmückt ist, so wie die \emph{Erde} mit
den Metallen, Pflanzen und Lebewesen\footnote{`animantibus'.} geschmückt
ist, die in ihr werden. Nachdem er über die Beschaffenheit des
Firmaments gesprochen hatte, blieb es ihm noch, über das Wirken der
Sterne\footnote{`siderum' und `gestirns'.} und ihre Funktion zu
sprechen, zu erklären, zu welchem \hspace{0pt}\hspace{0pt}Zweck sie
geschaffen und zu welcher Aufgabe sie von Gott bestimmt
wurden.\footnote{vom `Hexameron' (Ambrosius,
  \href{https://books.google.com/books?id=L1DbnAAACAAJ}{1897}, Buch 4,
  Kapitel 1).}

Die beiden Wirkungsweisen der Himmelskörper\footnote{`Celestium enim
  corporum' und `himlischen leiplichen ding'.}, \emph{Bewegung} und
\emph{Beleuchtung}, sind universell erkennbar, so ergibt sich eine
\emph{zweifache} Bewegung (I): Die eine ist (i) die gesamte
Welt\footnote{`mundi totius' und `ga{[}n{]}tze{[}n{]} werlt'.}, in der
sich Himmel und Äther in 24 Stunden in einem perfekten Kreis durch den
gesamten Raum des Universums\footnote{`vniuersi' und `werlt'.} drehen.
Die andere ist die eigentliche vielfältige und abwechslungsreiche (ii)
des Gestirns\footnote{`siderum' und `gestirns'.} mit der Hauptbewegung
der Sonne, die innerhalb von zwölf Monaten alle Tierkreiszeichen
durchläuft. Diese macht den Tag, daraus die Tagesabläufe, letztere
bildet das Jahr. Die übrigen Bewegungen der Sterne erfolgen in
unterschiedlichen Zeitabständen. Richtigerweise erinnerte uns
\emph{Moses} treffend und kurz daran, dass die Sterne für Tage, Jahre
und Zeiten in das Firmament gestellt seien.\footnote{vgl. `Heptaplus'
  (Mirandola \& Mirandola,
  \href{https://doi.org/10.3931/e-rara-61217}{1601}).}

Er wies auch ausdrücklich auf die übrige Wirkung der Sterne hin welches
die \emph{Beleuchtung}\footnote{`illuminatio' und `erlewchtu{[}n{]}g'.}
(II) ist, wenn er sagt, dass diese aufgestellt seien, um am Himmel zu
scheinen\footnote{`vt lucerent in celo' und `zu scheine{[}n{]} am{[}m{]}
  hymel'.} und die Erde zu erhellen. So werden die Körper des Mondes,
der Sonne und der Sterne\footnote{`stellarum' und `stern{[}n{]}'.} zu
diesem Zweck\footnote{`ministeria' und `dienstperkeiten'.}
verteilt\footnote{`distributa' und `außgetailt'.}. Die Sonne geht zwar
tagsüber nur \emph{scheinbar} alleinig auf (\textbf{Cicero} möchte wohl,
dass es so aussieht, dass nur \emph{eine} Sonne erkannt
wird)\footnote{bezieht sich auf \emph{De Natura Deorum} (Cicero \&
  Cicero, \href{https://doi.org/10.3931/e-rara-89116}{1496}, 2.68), wo
  angedeutet wird, dass die Sonne \emph{sol} (solus) genannt wird, weil
  sie die einzige ihrer Art ist, oder weil sie allein \emph{erscheint},
  nachdem sie alle anderen Sterne verdeckt hat, vgl. auch Lactantius
  (\href{https://books.google.com/books?id=qGsnghsi3uUC}{1497}, Buch 2,
  Kapitel 9/10).}, da sie durch die verhüllten Sterne\footnote{`sideribus'.}
einzig erscheint, dennoch ist \emph{sie} das wahre und vollkommene
Licht, das alles mit größter \emph{Wärme} und dem hellsten
\emph{Glanz}\footnote{`fulgore clarissimo' und `allerklarste{[}m{]}
  schein'.} erstrahlen lässt. Denn obwohl unzählige Sterne zu funkeln
und zu leuchten scheinen, spenden jene, deren Licht nicht voll und klar
ist, weder Wärme, noch vertreiben sie die Dunkelheit durch ihre
Vielzahl. So lassen sich \emph{zwei} Hauptelemente finden, (i)
\emph{Wärme} und (ii) \emph{`Glorie'}\footnote{`humor' und
  `feüchtigkeit' (um die Dunkelheit zu vertreiben).}, die eine
vielfältige und zugleich zusammenwirkende Kraft besitzen, von Gott auf
wunderbare Weise zur Erhaltung und Entstehung aller Dinge
geschaffen.\footnote{vom `Speculum Naturale' (De Beauvais,
  \href{https://nbn-resolving.org/urn:nbn:de:hbz:061:1-20612}{1964},
  Buch 1, Kapitel 75).}

Hier müssten allerdings \emph{sehr tiefgründige} Fragen behandelt
werden, von denen jede ein eigenes Buch erfordern würde: (i) In welcher
Weise befinden sich diese Sterne im Firmament, ob als dessen
\emph{edlere} Teile oder als Tiere\footnote{`a{[}n{]}{[}m{]}alia'(sic)
  und `geschöpff'.} in ihren Sphären, wie Fische im Wasser oder Tiere
auf der Erde.\footnote{spiegelt eine klassische philosophische Debatte
  des Mittelalters wider, sind Himmelskörper lediglich physische
  Bestandteile des Himmels oder verfügten sie über eine eigene
  \emph{Seele}.} Dieser Stelle erforderte ebenfalls eine
Auseinandersetzung mit den (ii) Astrologen\footnote{`Genethliacis'.}
hinsichtlich der Weissagung anhand der Sterne und der Vorhersage
zukünftiger Ereignisse, was die Wissenschaft darin bestätigte, dass
\emph{Moses} gesagt habe, die Sterne seien von Gott als Zeichen
platziert worden. Darüberhinaus muss (iii) die Natur der Sterne, ihre
Bewegung, ihre Umlaufbahnen, die \emph{Flecken}\footnote{`maculis'.} auf
dem Mond und die gesamte Wissenschaft der Sterne (Astronomie \emph{und}
Astrologie)\footnote{`siderali scie{[}n{]}tia'.} behandelt
werden.\footnote{aus dem `Heptaplus' (Mirandola \& Mirandola,
  \href{https://doi.org/10.3931/e-rara-61217}{1601}).}

Wahrlich,

\begin{quote}
qua{[}m{]}q{[}uoniam{]} sint pulcra {[}et{]} digna cognitu.\footnote{`Obwohl
  sie schön sind und Anerkennung verdienen' wird häufig verwendet um die
  Werke von Horaz
  (\href{https://google.cat/books?id=A28VAAAAYAAJ}{1868}) zu
  beschreiben, fasst sein berühmtes Prinzip aus der \emph{Ars Poetica}
  zusammen.}
\end{quote}

Vielleicht werden wir \emph{diesen} \textbf{Horaz} noch hören, aber
jetzt war nicht der richtige Ort.\footnote{`sed nunc non erat hic
  locus', Horace (\href{https://google.cat/books?id=A28VAAAAYAAJ}{1868},
  \emph{Ars Poetica}, Zeilen 14--19).}

\hypertarget{references}{%
\subsection{References}\label{references}}

Ambrosius, A. (1897). \emph{Opera: Exameron. De Paradiso. De Cain Et
Abel. De Noe. De Abraham. De Isaac. De Bono Mortis}. Edited by Faller,
O., \& Schenkl, K. Corpus Scriptorum Ecclesiasticorum Latinorum.
Hoelder-Pichler-Tempsky.
\url{https://books.google.com/books?id=L1DbnAAACAAJ}

Cicero, M. T., \& Cicero, Q. (1496). \emph{M.T. Ciceronis de Natura
Deorum Libri Tres} \ldots{} Impræssum Venetiis: per Symonem Papiensem
dictum Biuilaqua. \url{https://doi.org/10.3931/e-rara-89116}

De Beauvais, V. (1964). \emph{Speculum maius}. Graz: Akad. Druck- u.
Verl.-Anst. \url{https://nbn-resolving.org/urn:nbn:de:hbz:061:1-20612}

Horace. (1868). \emph{Q. Horatii Flacci Opera Omnia}. Oxoni et Londini:
J. Parker et Soc. \url{https://google.cat/books?id=A28VAAAAYAAJ}

Lactantius, L.C.F. (1497). \emph{Divinae Institutiones}. Venetiis: Simon
Bevilaqua. \url{https://books.google.com/books?id=qGsnghsi3uUC}

Mirandola, G., \& Mirandola, G. F. (1601). \emph{Ioannis Pici,
Mirandulae \ldots{} opera quae extant omnia \ldots{}} Basileae: per
Sebastianum Henricpetri. \url{https://doi.org/10.3931/e-rara-61217}

\end{document}
